% Este é o modelo do abnTeX2
% ---
% inserir lista de abreviaturas e siglas
% ---
%\begin{siglas}
%  \item[ABNT] Associação Brasileira de Normas Técnicas
%  \item[abnTeX] ABsurdas Normas para TeX
%\end{siglas}
% ---

%
% Outra opção é utilizar o pacote Acronym (mais funcional).
%

%\ac{SW}  	SW documents are beautifully typeset.
%\acf{SW} 	Scientific Word (SW) documents are beautifully typeset.
%\acs{SW} 	SW documents are beautifully typeset.
%\acl{SW} 	Scientific Word documents are beautifully typeset.

\pdfbookmark[0]{\listadesiglasname}{los}
\chapter*{\listadesiglasname}%

% Incluir no sumário
% \addcontentsline{toc}{chapter}{\listadesiglasname}%
\begin{acronym}
  \acro{AFD}{Autômato Finito Determinístico}
  \acro{AFN}{Autômato Finito não Determinístico}
  \acro{AFNL}{Autômato Finito não Determinístico Com Transições $\lambda$}
  \acro{ER}{Expressão Regular}
  \acro{APD}{Autômato de Pilha Determinístico}
  \acro{APN}{Autômato de Pilha não Determinístico}
  \acro{LLC}{Linguagem Livre de Contexto}
  \acro{GLC}{Gramática Livre de Contexto}
  \acro{MT}{Máquina de Turing}
  \acro{jFAST}{Java Finite Automata Simulation Tool}
  \acro{FSMD}{Finite State Machine Designer}
  \acro{GAM}{Ginux Abstract Machine}
\end{acronym}
